\chapter{Reger, Max}
\label{cha:reger}

\href{https://en.wikipedia.org/wiki/List_of_compositions_by_Max_Reger}{List of works} on Wikipedia.

\section{Orchestral}

\noindent{Zwei Romanzen, violin and small orchestra (1900), Op.~50, MR.~051 \\
\noindent{Sinfonietta, A major, two pianos (1904–05), Op.~90, MR.~107} \textit{Bamberger Symphoniker, Horst Stein}\marginnote{high} \\
\noindent{Serenade, G major (1905–06), Op.~95, MR.~112} \\
\noindent{Variations and Fugue on a Theme by Hiller, E major (1907), Op.~100, MR.~117} \\
\noindent{Violin Concerto, A major (1907–08), Op.~101, MR.~118} \textit{Ulf Wallinn, M\"{u}nchner Rundfunkorchester, Ulf Schirmer}\marginnote{high} \\
\noindent{Suite in A minor (vn, orch) (1908), Op.~103a, MR.~121} \textit{Walter Forchert, Bamberger Symphoniker, Horst Stein}\marginnote{high} \\
\noindent{Symphonischer Prolog zu einer Tragödie (1908), Op.~108, MR.~128} \\
\noindent{Piano Concerto, F minor (1910), Op.~114, MR.~139} \textit{Gerhard Oppitz, Bamberger Symphoniker, Horst Stein}\marginnote{high} \\
\noindent{Eine Lustspielouvertüre (1911), Op.~120, MR.~145} \textit{Bamberger Symphoniker, Horst Stein}\marginnote{high} \\
\noindent{Konzert im alten Stil, F major (1912), Op.~123, MR.~148} \textit{Bamberger Symphoniker, Horst Stein}\marginnote{high} \\
\noindent{Eine romantische Suite (1912), Op.~125, MR.~150} \\
\noindent{Vier Tondichtungen nach A. Böcklin (1913), Op.~128, MR.~153} \\
\noindent{Eine Ballett-Suite, D major (1913), Op.~130, MR.~155} \\
\noindent{Variations and Fugue on a Theme by Mozart (1914), Op.~132, MR.~160} \\
\noindent{Eine vaterländische Ouvertüre (1915), Op.~140, MR.~170} \\
\noindent{Andante and Rondo, (vn, orch) (1916), Op.~147} \textit{Bamberger Symphoniker, Horst Stein}\marginnote{high} \\
\noindent{Scherzino in C major (hn, str orch) (1899), WoO I/6, MR.~183} \textit{Marie-Luise Neunecker, Bamberger Symphoniker, Horst Stein}\marginnote{high} \\

\section{Chamber}

\noindent{Violin Sonata No. 1 in D minor (1890), Op.~1, MR.~001} \\
\noindent{Piano Trio No. 1 in B minor (vn, va, pf) (1891), Op.~2, MR.~002} \\
\noindent{Violin Sonata No. 2 in D major (1891), Op.~3, MR.~003} \\
\noindent{Cello Sonata No. 1 in F minor (1892), Op.~5, MR.~005} \\
\noindent{Cello Sonata No. 2 in G minor (1898), Op.~28, MR.~029} \\
\noindent{Violin Sonata No. 3 in A major (1899), Op.~41, MR.~042} \\
\noindent{Vier Sonaten (vn) (1900), Op.~42, MR.~043} \\
\noindent{Zwei Klarinettensonaten (cl/va, pf) (1900), Op.~49, MR.~050} \\
\noindent{String Quartet No. 1 in G minor (1901), Op.~54/1} \\
\noindent{String Quartet No. 2 in A major (1901), Op.~54/2} \\
\noindent{Piano Quintet No. 2 in C minor (1901–1902), Op.~64, MR.~064} \\
\noindent{Violin Sonata No. 4 in C major (1903), Op.~72, MR.~072} \\
\noindent{String Quartet No. 3 in D minor (1903–1904), Op.~74, MR.~074} \\
\noindent{Serenade No. 1 in D major (fl, vn, va) (1904), Op.~77a, MR.~082} \\
\noindent{String Trio No. 1 in A minor (1904), Op.~77b, MR.~083} \\
\noindent{Cello Sonata No. 3 in F major (1904), Op.~78, MR.~084} \\
\noindent{Kompositionen (vn, pf) (1902–1904), Op.~79d, MR.~089} \\
\noindent{Kompositionen (vc, pf) (1902–1904), Op.~79e, MR.~090} \\
\noindent{Violin Sonata No. 5 in F-sharp minor (1905), Op.~84, MR.~100} \\
\noindent{Zwei Kompositionen (vn, pf) (1905), Op.~87, MR.~104} \\
\noindent{Seven Violin Sonatas (vn) (1905), Op.~91, MR.~108} \\
\noindent{Suite im alten Stil in F major (vn, pf) (1906), Op.~93, MR.~110} \\
\noindent{Piano Trio No. 2 in E minor (vn, vc, pf) (1907–1908), Op.~102, MR.~119} \\
\noindent{Hausmusik (vn, pf) (1908–1909), Op.~103, MR.~120} \\
\noindent{Suite in A minor (vn, pf) (1908), Op.~103a, MR.~121} \\
\noindent{Violin Sonata No. 6 in D minor (1909), Op.~103b/1, MR.~122} \\
\noindent{Violin Sonata No. 7 in A major (1909), Op.~103b/2, MR.~122[1]} \\
\noindent{Zwölf kleine Stücke nach eigenen Liedern (vn, pf) (1909), Op.~103c, MR.~123} \\
\noindent{Clarinet Sonata No. 3 in B-flat major (cl/va, pf) (1908–1909), Op.~107, MR.~127} \\
\noindent{String Quartet No. 4 in E-flat major (1909), Op.~109, MR.~129} \\
\noindent{Piano Quartet No. 1 in D minor (1910), Op.~113, MR.~138} \\
\noindent{Cello Sonata No. 4 in A minor (1910), Op.~116, MR.~141} \\
\noindent{Präludien und Fugen für Violine solo (vn) (1909–1912), Op.~117, MR.~142} \\
\noindent{String Sextet in F major (1910), Op.~118, MR.~143} \\
\noindent{String Quartet No. 5 in F-sharp minor (1911), Op.~121, MR.~146} \\
\noindent{Violin Sonata No. 8 in E minor (1911), Op.~122, MR.~147} \\
\noindent{Präludien und Fugen (vn) (1914), Op.~131a, MR.~156} \\
\noindent{Three Duos (canons and fugues) in old style (2 vn) (1914), Op.~131b, MR.~157} \\
\noindent{Drei Suiten (vc) (1915), Op.~131c, MR.~158} \\
\noindent{Drei Suiten (va) (1915), Op.~131d, MR.~159[2]} \\
\noindent{Piano Quartet No. 2 in A minor (1914), Op.~133, MR.~162} \\
\noindent{Violin Sonata No. 9 in C minor (1915), Op.~139, MR.~169} \\
\noindent{Serenade No. 2 in G major (fl, vn, va) (1915), Op.~141a, MR.~171} \\
\noindent{String Trio No. 2 in D minor (1915), Op.~141b, MR.~172} \\
\noindent{Clarinet Quintet in A major (1915–1916), Op.~146, MR.~179} \textit{Kilian Herold, Armida Quartett}\marginnote{01/2024, max} \\
\noindent{Scherzino in C major (hn, str) (1899), WoO I/6, MR.~183} \\
\noindent{Wind Serenade (first movement) in B-flat major (1904), WoO I/9} \\
\noindent{Scherzo in G minor (fl, str qt) (1888 or 1889), WoO II/1} \\
\noindent{String Quartet No. 0 in D minor (str qt, db) (1888–1889), WoO II/2, MR.~184} \\
\noindent{Scherzo in G minor (2 str qt) (1890–1892), WoO II/6} \\
\noindent{Piano Quintet No. 1 in C minor (1897–1898), WoO II/9, MR.~185} \\
\noindent{Romanze in G major (vn, pf) (1901), WoO II/10, MR.~187} \\
\noindent{Petite Caprice in G minor (vn, pf) (1901), WoO II/11, MR.~188} \\
\noindent{Tarantella in G minor (cl/vn, pf) (1901–1902), WoO II/12, MR.~188} \\
\noindent{Albumblatt in E-flat major (cl/vn, pf) (1902), WoO II/13, MR.~195} \\
\noindent{Allegretto grazioso in A major (fl, pf) (1902), WoO II/14, MR.~194} \\
\noindent{Caprice in A minor (vc, pf) (1902), WoO II/15, MR.~193} \\
\noindent{Prelude and Fugue in A minor (vn) (1902), WoO II/16, MR.~190} \\
\noindent{Allegro in A major (2 vn) (1907?), WoO II/18, MR.~189} \\
\noindent{Prelude in E minor (vn) (1915), WoO II/19, MR.~192} \\


\section{Piano}

\noindent{Walzercapricen, 4-hand piano (1892), MR.~9} \\
\noindent{Deutsche Tänze, 4-hand piano (1892), MR.~10} \\
\noindent{Sieben Walzer (1893), MR.~11} \\
\noindent{Lose Blätter, Kleine Stücke  (1894), MR.~13} \\
\noindent{Aus der Jugendzeit, Zwanzig kleine Stücke (1895), MR.~18} \\
\noindent{Improvisationen (1896), MR.~19} \\
\noindent{Fünf Humoresken (1898), MR.~21} \\
\noindent{Sechs Walzer, 4-hand piano (1898), MR.~23} \\
\noindent{Six Morceaux (1898), MR.~25} \\
\noindent{Aquarellen, Kleine Tonbilder (1897–98), MR.~26} \\
\noindent{Sieben Fantasiestücke (1898), MR.~27} \\
\noindent{Sieben Charakterstücke (1899), MR.~33} \\
\noindent{Cinq Pièces pittoresques, 4-hand piano (1899), MR.~35} \\
\noindent{Bunte Blätter, Kleine Stücke (1899), MR.~37} \\
\noindent{Zehn kleine Vortragsstücke, zum Gebrauch beim Unterricht (1900), MR.~45} \\
\noindent{Sechs Intermezzi (1900), MR.~46} \\
\noindent{Silhouetten, Sieben Stücke (1900), MR.~54} \\
\noindent{Sechs Burlesken, 4-hand piano (1900), MR.~58} \\
\noindent{Haus- und Kirchenmusik (1900–04), MR.~85} \\
\noindent{Kompositionen (1900–04), MR.~86} \\
\noindent{Variationen und Fuge über ein Thema von Joh. Seb. Bach (1904), MR.~94} \textit{Eden Walker}\marginnote{07/2024, max} \\
\noindent{Aus meinem Tagebuche (Four volumes) (1904–12), MR.~95} \\
\noindent{Variationen und Fuge über ein Thema von Beethoven, two pianos (1904), MR.~102} \\
\noindent{Vier Sonatinen, two pianos (1905–08), MR.~106} \\
\noindent{Sechs Stücke, 4-hand piano (1906), MR.~111} \\
\noindent{Introduction, Passacaglia and Fugue, two pianos (1906), MR.~113} \\
\noindent{Six Preludes and Fugues, piano (1907), MR.~116} \\
\noindent{Acht Episoden, Klavierstücke für grosse und kleine Leute (1910), MR.~140} \\
\noindent{Variations and Fugue on a Theme by Telemann (1914), MR.~163} \\
\noindent{Träume am Kamin, Zwölf kleine Stücke (1915), MR.~174} \textit{Eden Walker}\marginnote{07/2024, max} \\
\noindent{Canons in all major and minor tonalities (1895), MR.~196} \\
\noindent{Grüße an die Jugend (1898), MR.~198} \\
\noindent{Liebestraum (1898), MR.~182} \\
\noindent{Miniature Gavotte (1898), MR.~199} \\
\noindent{Improvisation über den Walzer von Johann Strauß "An der schönen blauen Donau" (1898), MR.~200} \\
\noindent{Blätter und Blüten (1900–02), MR.~202} \\
\noindent{Vier Spezialstudien für die linke Hand (1901), MR.~204} \\
\noindent{In der Nacht (1902), MR.~205} \\
\noindent{Perpetuum mobile (1905), MR.~206} \\
\noindent{Scherzo (1906), MR.~207} \\
\noindent{Caprice (1906), MR.~208} \\
\noindent{Ewig Dein! (1907), MR.~181} \\
\noindent{Fughetta on the "Deutschlandlied" (1916), MR.~209}
